\documentclass[12pt]{exam}
\usepackage[phy]{template-for-exam}
\usepackage{empheq}

\newcommand{\printeqs}{
  \begin{center}
    \vspace{-1cm}
    \begin{empheq}[box=\fbox]{align*}
     &
     & 
      v &= \frac{d}{t}        &    
      a &= \frac{\Delta v}{t} &
      \Delta v &= v_f - v_i   &
     &
    \end{empheq}
  \end{center}
}

\title{Motion \#4}
\author{Rohrbach}


\begin{document}
\maketitle

\printeqs

\noindent
{\bf Please use the proper problem-solving method.}  (1) Draw a picture; (2) knowns \& unknowns; (3) pick an equation; (4) plug \& chug; (5) answer with units.


\begin{questions}


\question
  What is the acceleration of a unicycle that goes from 
  3 m/s to 5 m/s in 7 s?
  \vs

\question 
  A train is traveling at \SI{32}{\meter\per\second} and begins slowing down at \SI{-4}{\meter\per\second^2}. What will its velocity be after \SI{5}{\second}? 
  \vs

\question
  A hiker walks at a constant velocity of \SI{1.5}{\meter\per\second}. How long will it take the hiker to travel {\bf 1.6 km}? 
  \vs

\pagebreak

\printeqs

\question
  What is the velocity of an ATV that travels 13 m in 2 s?
  \vs

\question
  A car is traveling at \SI{24}{\meter\per\second}. How far will it travel in {\bf 2.5 minutes} if it maintains this velocity?
  \vs

\question
  A skateboarder is traveling at \SI{6}{\meter\per\second} and accelerates at a rate of \SI{2.5}{\meter\per\second^2} for \SI{4.0}{\second}. What is the skateboarder's final velocity?
  \vs

  
\end{questions}
\end{document}